\subsection{Distributional Evolution via Prompt Engineering}
\label{subsec:resultados_prompts}

Table~\ref{tab:resultados_prompts} presents the evolution of the
Jensen-Shannon Divergence (JSD) across the four experimental scenarios for
question P02 (Policy Priorities), utilizing the Sabiazinho-4 model.

\begin{table}[htbp]
    \centering
    \caption{JSD evolution across experimental scenarios (P02~--- Policy
        Priorities, Sabiazinho-4 model).}
    \label{tab:resultados_prompts}
    \begin{tabular}{@{}llcc@{}}
        \toprule
        \textbf{Scenario} & \textbf{Applied Technique} & \textbf{JSD} & \textbf{$\Delta$JSD} \\
        \midrule
        \textit{Baseline} & Raw \textit{prompt} (13 attributes) & 0.525 & --- \\
        Scenario 1 & Feature selection (7 attr.) & 0.366 & $-0.159$ \\
        Scenario 2 & Natural prose \textit{prompt} & 0.366 & $0.000$ \\
        Scenario 3 & \textit{System/User} + CoT & 0.297 & $-0.069$ \\
        Scenario 4 & \textit{Few-shot} (3 personas) & 0.121 & $\mathbf{-0.176}$ \\
        \bottomrule
    \end{tabular}
\end{table}

We observed an overall reduction of 77\% in JSD between the
\textit{baseline} ($0.525$) and the best experimental scenario ($0.121$), as
illustrated in Figure~\ref{fig:evolucao_jsd}. Feature selection (Scenario~1)
produced an immediate and substantial improvement ($\Delta\text{JSD} = -0.159$),
indicating that eliminating informational noise helps the model anchor its
decisions on the most decisive variables. Natural prose reformatting
(Scenario~2) maintained the identical JSD value, suggesting that textual syntax
alone does not alter the output probability distribution. Conversely, the Chain
of Thought reasoning (Scenario~3) promoted a significant incremental gain
($\Delta\text{JSD} = -0.069$). Finally, introducing calibration examples
(\textit{few-shot}, Scenario~4) allowed the simulation to surpass the target
threshold of distributional convergence ($\text{JSD} < 0.15$), reaching
$0.121$.

\begin{figure}[htbp]
    \centering
    \resizebox{\columnwidth}{!}{%
        \begin{tikzpicture}[
    font=\sffamily\scriptsize,
    xscale=1.5,
    yscale=3.0,
    bar/.style={draw=black!85, line width=0.5pt},
    refline/.style={dashed, draw=black!60, line width=0.5pt},
    label/.style={font=\tiny\bfseries}
]

    % --- Y-axis grid ---
    \foreach \y in {0.0, 0.1, 0.2, 0.3, 0.4, 0.5, 0.6} {
        \draw[black!15, line width=0.4pt] (-0.3, \y) -- (4.5, \y);
        \node[left, font=\tiny, text=black!75] at (-0.3, \y) {\pgfmathprintnumber[fixed,precision=1]{\y}};
    }
    \draw[black!85, line width=0.6pt] (-0.3,0) -- (4.5,0);
    \draw[black!85, line width=0.6pt] (-0.3,0) -- (-0.3,0.65);
    \node[rotate=90, font=\scriptsize\bfseries, text=black!85] at (-0.75, 0.30) {JSD ($\downarrow$ better)};

    % --- Reference line ---
    \draw[refline] (-0.3, 0.15) -- (4.5, 0.15);
    \node[right, font=\tiny\itshape, text=black!70] at (4.5, 0.15) {JSD = 0.15};

    % --- Scenario bars ---
    \filldraw[bar, fill=black!75] (-0.25, 0) rectangle (0.25, 0.525);
    \node[above, label] at (0, 0.525) {0.525};
    \node[below, font=\tiny, align=center] at (0, -0.02) {Baseline\\(13 attrs.)};

    \filldraw[bar, fill=black!55] (0.75, 0) rectangle (1.25, 0.366);
    \node[above, label] at (1.0, 0.366) {0.366};
    \node[font=\tiny\bfseries, text=white] at (1.0, 0.183) {$\Delta -0.159$};
    \node[below, font=\tiny, align=center] at (1.0, -0.02) {C1: feature\\selection\\(7 attrs.)};

    \filldraw[bar, fill=black!55] (1.75, 0) rectangle (2.25, 0.366);
    \node[above, label] at (2.0, 0.366) {0.366};
    \node[below, font=\tiny, align=center] at (2.0, -0.02) {C2: natural\\prose};

    \filldraw[bar, fill=black!35] (2.75, 0) rectangle (3.25, 0.297);
    \node[above, label] at (3.0, 0.297) {0.297};
    \node[font=\tiny\bfseries, text=white] at (3.0, 0.148) {$\Delta -0.069$};
    \node[below, font=\tiny, align=center] at (3.0, -0.02) {C3: system/user\\+ CoT};

    \filldraw[bar, fill=black!15] (3.75, 0) rectangle (4.25, 0.121);
    \node[above, label, fill=white, inner sep=0.8pt] at (4.0, 0.121) {0.121};
    \node[font=\tiny\bfseries, text=black!90] at (4.0, 0.055) {$\Delta -0.176$};
    \node[below, font=\tiny, align=center] at (4.0, -0.02) {C4: few-shot\\(3 personas)};

    % --- Global reduction ---
    \draw[->, >=Stealth, line width=0.8pt, draw=black!90] (0.0, 0.55) to[out=18, in=162] node[above=1pt, font=\scriptsize\bfseries, fill=white, inner sep=1pt] {$-77\%$} (4.0, 0.15);
\end{tikzpicture}
%
    }
    \caption{JSD evolution across prompt engineering scenarios (P02,
        Sabiazinho-4), demonstrating a 77\% reduction from baseline to
        Scenario~4 (few-shot).}
    \label{fig:evolucao_jsd}
\end{figure}

\subsection{Systematic Benchmark: Brazilian Portuguese-Centric vs. Global Models}
\label{subsec:resultados_modelos}

Table~\ref{tab:resultados_modelos} summarizes the comparative performance of
the four evaluated models (Sabiá-4, Sabiazinho-4, GPT-5-mini, and
LLaMA~3.2~3B) across the three target questions (P02, P03, and P04).

\begin{table}[htbp]
    \centering
    \begin{threeparttable}
        \caption{Comparative performance of models across questions using JSD.}
        \label{tab:resultados_modelos}
        \setlength{\tabcolsep}{6pt}
        \begin{tabular}{@{}llc@{}}
            \toprule
            \textbf{Question} & \textbf{Model} & \textbf{JSD ($\downarrow$)} \\
            \midrule
            \multirow{4}{*}{\shortstack[l]{P02: Policy\\Priorities}}
                & Sabiá-4 & 0.7238 \\
                & Sabiazinho-4 & 0.5664 \\
                & GPT-5-mini & 0.6185 \\
                & LLaMA~3.2~3B & \textbf{0.4882} \\
            \midrule
            \multirow{4}{*}{\shortstack[l]{P03: Fake\\News}}
                & Sabiá-4 & 0.6846 \\
                & Sabiazinho-4 & \textbf{0.5452} \\
                & GPT-5-mini & 0.5908 \\
                & LLaMA~3.2~3B & 0.5912 \\
            \midrule
            \multirow{4}{*}{\shortstack[l]{P04: Political\\Participation}}
                & Sabiá-4 & 0.3007 \\
                & Sabiazinho-4 & \textbf{0.2097} \\
                & GPT-5-mini & 0.8263 \\
                & LLaMA~3.2~3B & 0.8100 \\
            \midrule
            \multirow{4}{*}{\shortstack[l]{Mean JSD\\($\overline{\text{JSD}}$)}}
                & Sabiazinho-4 & \textbf{0.4404} \\
                & Sabiá-4 & 0.5697 \\
                & LLaMA~3.2~3B & 0.6298 \\
                & GPT-5-mini & 0.6785 \\
            \midrule
            \multicolumn{2}{@{}l}{\textbf{Overall Mean $\pm$ SD}} & \textbf{0.5796 $\pm$ 0.1839} \\
            \bottomrule
        \end{tabular}
        \begin{tablenotes}
            \footnotesize
            \item \textit{Note:} $\downarrow$ indicates that lower values denote
                closer alignment to the real CESOP distribution (JSD). Bold indicates best performance per question and overall mean.
        \end{tablenotes}
    \end{threeparttable}
\end{table}

The aggregated comparative analysis reveals that Sabiazinho-4 achieved
the best average distributional adherence ($\overline{\text{JSD}} = 0.4404$),
followed by Sabiá-4 ($\overline{\text{JSD}} = 0.5697$),
LLaMA~3.2~3B ($\overline{\text{JSD}} = 0.6298$), and
GPT-5-mini ($\overline{\text{JSD}} = 0.6785$).

Analyzing results across individual questions (Table~\ref{tab:resultados_modelos}) demonstrates
that simulation accuracy is strongly conditioned by thematic domain characteristics.
On the one hand, Question P04 (Political Participation) revealed the most pronounced
polarization between model families: Brazilian architectures exhibited high empirical
alignment ($\text{JSD} = 0.2097$ for Sabiazinho-4 and $0.3007$ for Sabiá-4), whereas
global baselines failed substantially ($\text{JSD} = 0.8263$ for GPT-5-mini and
$0.8100$ for LLaMA~3.2~3B) due to acute social desirability bias. On the other hand,
Questions P02 and P03 (Policy Priorities and Fake News) showed moderate dispersion
across all models ($\text{JSD}$ between $0.48$ and $0.72$), reflecting their larger
combinatorial choice spaces (14 and 8 alternatives) and underlying socioeconomic
complexity.

Regarding cross-consistency across questions and models, Sabiazinho-4
maintained superior stability across all thematic dimensions, avoiding the
severe degradation peaks observed in globally aligned models
(Table~\ref{tab:resultados_modelos}).

\subsection{Sociocultural Biases, Explainability Auditing, and Methodological Trade-offs}
\label{subsec:discussao_qualitativa}

Detailed inspection of simulated distributions revealed systematic bias
patterns conditioned by model origin and training alignment. Brazilian
architectures (Sabiá-4 and Sabiazinho-4) demonstrated appropriate sensitivity to
the civic skepticism of the Brazilian electorate in P04, accurately capturing
the substantial proportion of respondents expressing little or no desire to
participate in municipal politics. However, in P02, they exhibited a tendency to
over-prioritize healthcare quality, occasionally assigning it to higher-income
personas holding private health insurance. Conversely, global models (GPT-5-mini
and LLaMA~3.2~3B) displayed pronounced social desirability and institutional
optimism biases: in P04, GPT-5-mini attributed positive participation desire to
over 98\% of synthetic personas, failing to reproduce documented political
alienation, while responses in P02 heavily concentrated on job creation in a
standard corporate/institutional style.

Incorporating intermediate reasoning (Chain-of-Thought) in Scenario~3 enabled
qualitative auditing of persona decision processes, resolving the opacity
inherent in unreasoned prompts. Qualitative inspection of explanatory rationales
unveiled two primary phenomena. First, genuine demographic alignment occurred
when models anchored decisions at the intersection of income, education, and
geography (for instance, lower-income, less-educated personas articulating civic
disengagement due to perceived institutional corruption and distance). Second,
persona dissonance emerged when models conflated simulated profiles with generic
knowledge (such as affluent personas holding private coverage selecting public
healthcare improvements due to broader societal perceptions of health system
precarity). This transparency establishes XAI not merely as a \textit{post hoc}
validation tool, but as an active diagnostic instrument for auditing behavioral
coherence in synthetic populations.

Finally, while Scenario~4 attained optimal distributional convergence ($\text{JSD} =
0.121$), introducing manual calibration examples entails an essential
methodological trade-off: in-context demonstrations serve as potent
distributional anchors, requiring balanced selection to prevent steering models
toward researcher priors.
